\documentclass[10pt]{article}
\usepackage[T1]{fontenc}
\usepackage[utf8]{inputenc}
\usepackage[french]{babel}
\usepackage{amsmath,amssymb,mathrsfs,mathpazo}
\usepackage{geometry}
\geometry{a4paper,margin=2cm}
\usepackage{graphicx,array,tabularx,tabularray,colortbl,multirow,multicol,booktabs,enumitem,pifont,ulem,fancybox,diagbox,cancel}
\usepackage{pstricks,pstricks-add,pst-plot,pst-tree,pst-node,pst-text,pst-func,pst-math,pst-bspline,pst-3dplot,pst-eucl,pst-solides3d}
\let\pscalculate\undefined
\usepackage{pst-bezier}
\usepackage{tikz,pgfplots,tkz-tab}
\usetikzlibrary{arrows,arrows.meta,patterns,calc,positioning,trees,decorations.markings,shapes.geometric,shapes.misc,shapes.arrows,intersections,decorations.pathmorphing,decorations.pathreplacing,backgrounds,shadows,babel}
\pgfplotsset{compat=1.18}
\usepgfplotslibrary{fillbetween}
\usepackage[np]{numprint}
\usepackage{listings,makecell,pas-tableur,twoopt}
\let\barre\undefined
\usepackage[pstricks]{tablvar}
\definecolor{darkgreen}{RGB}{0,100,0}
\providecommand{\red}{\color{red}}
\providecommand{\blue}{\color{blue}}
\providecommand{\green}{\color{green}}
\providecommand{\black}{\color{black}}
\providecommand{\white}{\color{white}}
\providecommand{\gray}{\color{gray}}
\providecommand{\cyan}{\color{cyan}}
\providecommand{\magenta}{\color{magenta}}
\providecommand{\yellow}{\color{yellow}}
\providecommand{\violet}{\color{violet}}
\providecommand{\orange}{\color{orange}}
\providecommand{\euro}{\texteuro}
\providecommand{\EUR}{\texteuro}
\providecommand{\ieme}{\textsuperscript{e}}
\providecommand{\degres}{\ensuremath{^\circ}}
\providecommand{\degree}{\ensuremath{^\circ}}
\providecommand{\vv}[1]{\overrightarrow{#1}}
\providecommand{\Coord}[1]{#1}
\providecommand{\decosix}{}
\providecommand{\starredbullet}{\ensuremath{\bullet}}
\pagestyle{empty}
\setlength{\parindent}{0pt}
\usepackage{adjustbox}\DeclareUnicodeCharacter{2236}{:}\DeclareUnicodeCharacter{2190}{\ensuremath{\leftarrow}}
\providecommand{\R}{\mathbb{R}}
\providecommand{\N}{\mathbb{N}}
\providecommand{\Z}{\mathbb{Z}}
\providecommand{\Q}{\mathbb{Q}}
\providecommand{\D}{\mathbb{D}}
\providecommand{\C}{\mathbb{C}}
\providecommand{\e}{{\mathrm{e}}}
\providecommand{\vect}[1]{\overrightarrow{#1}}
\providecommand{\vectt}[1]{\overrightarrow{\mathrm{#1}}}
\providecommand{\vvt}[1]{\overrightarrow{\mathrm{#1}}}
\providecommand{\np}[2][]{{#2~#1}}
\providecommand{\vv}[1]{{\overrightarrow{#1}}}
\renewcommand{\d}{{\mathrm{d}}}
\providecommand{\degres}{{^\circ}}
\providecommand{\degre}{{^\circ}}
\providecommand{\blue}{\color{blue}}
\providecommand{\red}{\color{red}}
\providecommand{\enskip}{\quad}
\providecommand{\llbracket}{\lbrack\!\lbrack}
\providecommand{\rrbracket}{\rbrack\!\rbrack}
\providecommand{\euro}{{\text{€}}}
\providecommand{\no}{{\text{n°}}}
\providecommand{\interval}[2]{{[#1;#2]}}
\providecommand{\intervFO}[2]{{[#1;#2[}}
\providecommand{\dotfill}{\ldots}

\providecommand{\ds}{\displaystyle}
\providecommand{\textding}[1]{#1}
\providecommand{\Oij}{\(\mathrm{O};\vec{\imath},\vec{\jmath}\)}
\providecommand{\Oijk}{\(\mathrm{O};\vec{\imath},\vec{\jmath},\vec{k}\)}
\providecommand{\Ouv}{\(\mathrm{O};\vec{u},\vec{v}\)}
\providecommand{\intervFF}[2]{\ensuremath{[#1;#2]}}
\renewcommand{\euro}{\ifmmode\text{\texteuro}\else\texteuro\fi}
\DeclareUnicodeCharacter{00B0}{\ifmmode^{\circ}\else\textdegree\fi}
\DeclareUnicodeCharacter{20AC}{\ifmmode\text{\texteuro}\else\texteuro\fi}
% Énoncés : Éducation nationale. Transcription et corrigés : APMEP, Denis Vergès et contributeurs. Découpage et classement : CoopMaths, Sébastien Lozano, Éric Elter et contributeurs. Adaptation de la présentation : Maths Explorer.
% CC BY-SA : https://coopmaths.fr/www/about/
% https://coopmaths.fr/alea/static/bac/2023/tex/bac_2023_03_sujet2_asie_3_cor.tex
\begin{document}
\noindent\textit{APMEP / CoopMaths. CC BY-SA. Présentation adaptée par Maths Explorer.}\par\medskip
\medskip

On décide d'étudier le rayonnement radioactif du polonium lors de la désintégration des noyaux atomiques au cours du temps.

Au début de l'expérience, on dispose d'un morceau de $2$~g de polonium.

On sait que $1$~g de polonium contient $3 \times 10^{21}$ noyaux atomiques.

On admet que, au bout de 24 heures, $0,5\,\%$ des noyaux se sont désintégrés et que, pour compenser cette disparition, on ajoute alors $0,005$ g de polonium.

On modélise la situation à l'aide d'une suite $\left(v_n\right)_{n \in \N}$ ; on note $v_0$ le nombre de noyaux contenus dans le polonium au début de l'expérience. 

Pour $n \geqslant 1,\: v_n$ désigne le nombre de noyaux contenus dans le polonium au bout de $n$ jours écoulés.

\medskip

\begin{enumerate}
\item 
	\begin{enumerate}
		\item 
$v_0 = 2 \times 3 \times 10^{21} = 6 \times 10^{21}$noyaux atomiques.
		\item 
Si $v_n$ est le nombre de noyaux contenus dans le polonium au bout de $n$ jours, le lendemain 0,5\,\% ont disparu, il reste donc $v_n \left( 1 -\dfrac{0,5}{100}\right)  = v_n \times (1 - 0,005) = 0,995v_n$ que l'on augmente de 0,005. Donc $v_{n+1} = 0,995v_n + 0,005 \times 3 \times 10^{21} = 0,995v_n + 0,015 \times 10^{21}$ ou encore 

\[v_{n+1} = 0,995v_n + 1,5 \times 10^{19}.\]
	\end{enumerate}
\item 
	\begin{enumerate}
		\item 
\emph{Initialisation} :

Pour $n = 0,\:u_0 = 6\times 10^{21}$ ;

Pour $n = 1$, la relation de récurrence donne :

$u_1 = 0,995\times 6\times 10^{21}v + 1,5 \times 10^{19} = 5,97 \times 10^{21} + 1,5 \times 10^{19} = 597 \times 10^{19} + 1,5 \times 10^{19} = 598,5 \times 10^{19} = 5,985 \times 10^{21}$.

On a donc $0\leqslant v_{1}  \leqslant v_0$. la proposition est vraie au rang $n = 0$.

\emph{Hérédité} : soit $n \in \N$ tel que $0 \leqslant v_{n+1}  \leqslant v_n$. On peut écrire successivement :	

$0 \leqslant v_{n+1} \leqslant v_n \iff 0,995\times 0 \leqslant 0,995v_{n+1} \leqslant 0995v_n \iff 0 + 1,5 \times 10^{19} \leqslant v_{n+1} \times 0,995 + 1,5 \times 10^{19} \leqslant 0,995v_n + 1,5 \times 10^{19}$ soit $1,5 \times 10^{19} \leqslant u_{n+2} \leqslant u_{n+1}$. La proposition est donc vraie au rang $n + 1$.

\emph{Conclusion }: la proposition est vraie au rang 0 et si elle est vraie au rang $n$ elle l'est aussi au rang $n + 1$ : d'après le principe de récurrence, pour tout $n \in \N$ :
\[0 \leqslant v_{n+1}  \leqslant v_n.\]

		\item 
		
La partie gauche de la proposition montre que la suite $\left(v_n\right)_{n\in \N}$ est minorée par $0$ et la partie de droite montre que la suite $\left(v_n\right)_{n\in \N}$ est décroissante, donc la suite $\left(v_n\right)_{n\in \N}$ est convergente vers un réel $\ell \geqslant 0$.
	\end{enumerate}
\item On considère la suite $\left(u_n\right)_{n \in \N}$ définie, pour tout entier naturel $n$, par: \[u_n = v_n - 3 \times 10^{21}.\]
	\begin{enumerate}
		\item 
D'après la définition on a :

$u_{n+1} = v_{n+1} - 3 \times 10^{21} = 0,995v_n + 1,5 \times 10^{19} - 3 \times 10^{21} = 0,995v_n + 1,5 \times 10^{19} - 300 \times 10^{19}$, soit $u_{n+1} = 0,995v_n - 298,5 \times 10^{19} = 0,995v_n - 2,985 \times 10^{21} = 0,995v_n - 3 \times 0,995\times 10^{21} = 0,995\left(v_n - 3\times 10^{21} \right) = 0,995 u_n$.

L'égalité $u_{n+} = 0,995u_n$, vraie pour tout $n \in \N$ montre que la suite $\left(u_n\right)_{n \in \N}$ est géométrique de raison $0,995$, de premier terme $u_0 = v_0 - 3 \times 10^{21} = 6 \times 10^{21} - 3 \times 10^{21} = 3 \times 10^{21}$.
		\item 
On sait que le terme général de la suite $\left(u_n\right)_{n \in \N}$ est $u_n = u_0 \times 0,995^n = 3 \times 10^{21} \times 0,995^n$.

Or $u_n = v_n - 3 \times 10^{21} \iff v_n = u_n + 3 \times 10^{21} = 3 \times 10^{21} \times 0,995^n + 3 \times 10^{21} = 3 \times 10^{21}\left(0,995^n + 1\right)$.


		\item 
		Comme $- 1 < 0,995 < 1$, on sait que $\displaystyle\lim_{n \to + \infty} 0,995^n = 0$, donc $\displaystyle\lim_{n \to + \infty}0,995^n + 1 = 1$ et enfin $\displaystyle\lim_{n \to + \infty}v_n = 3 \times 10^{21}$.
	\end{enumerate}
\item 
Il faut résoudre l'inéquation :

$v_n < 4,5 \times 10^{21} \iff 3 \times 10^{21}\left(0,995^n + 1\right) < 4,5 \times 10^{21} \iff 3\left(0,995^n + 1\right) < 4,5 \iff 0,995^n + 1 < 1,5 \iff 0,995^n < 0,5 \iff n \ln 0,995 < \ln 0,5 \iff n > \dfrac{\ln 0,5}{\ln 0,995}$ (car $\ln 0,995 < 0$).

La calculatrice donne $\dfrac{\ln 0,5}{\ln 0,995} \approx 138,3$.

Le nombre de noyaux de polonium sera inférieur à $4,5 \times 10^{21}$ au bout de 139 jours.
\item 















	\begin{enumerate}
		\item 
Pour définir V on peut utiliser la définition ou la formule explicite, soit 

V = 0.995*V + 1.5*10**19 ou

V= 3*10**21*(0.995**(k+1)+1)
		\item 
52 semaines à 7 jours représentent 364 jours. Il faut donc écrire noyaux(364).
	\end{enumerate}
\end{enumerate}

\medskip

\end{document}
